\documentclass[../../main.tex]{subfiles}
\begin{document}

\begin{center}
  {\fontsize{48pt}{58pt}\selectfont \textbf{1994}\normalsize\textbf{年度}}
\end{center}

\begin{table}[htb]
  \centering
  \caption{出題テーマと難易度}
  \begin{tabular}{|c|p{9cm}|c|}
    \hline
    \textbf{問題番号} & \textbf{テーマ} & \textbf{難易度} \\
    \hline
    第1問 & 多項式の正値性と実根の個数 & \\
    \hline
    第2問 & 三角関数の値と整数になる数列 & \\
    \hline
    第3問 & 立体の断面積と体積 & \\
    \hline
    第4問 & 積分の漸化式と関数の決定 & \\
    \hline
    第5問 & カードの確率と期待値の最大化 & \\
    \hline
    第6問 & 格子距離による軌跡 & \\
    \hline
  \end{tabular}
\end{table}

\begin{table}[htb]
  \centering
  \caption{解答}
  \renewcommand{\arraystretch}{1.4}
  \begin{tabular}{|c|c|p{9cm}|}
    \hline
    \textbf{問題番号} & \textbf{小問} & \textbf{答え} \\
    \hline
    第1問 & (1) & $f(x)>0$（証明問題） \\
    \cline{2-3}
    & (2) & 実根は$-1<\alpha<0$にただ1つ \\
    \hline
    第2問 & (1) & $a+b=\dfrac{5}{4}$，$ab=\dfrac{5}{16}$ \\
    \cline{2-3}
    & (2) & 与えられた数はすべての自然数$n$で整数（証明問題） \\
    \hline
    第3問 & (1) & $S=\cos^2\theta\left(\theta-\dfrac{1}{2}\sin 2\theta\right)$ \\
    \cline{2-3}
    & (2) & $V=\dfrac{4}{45}$ \\
    \hline
    第4問 & -- & $f(x)=\dfrac{1}{(1-x)^2}$ \\
    \hline
    第5問 & (1) & $E=-p^2+p-2q^2+4q$ \\
    \cline{2-3}
    & (2) & 最大値は$p=\dfrac{1}{6}$，$q=\dfrac{1}{5}$のとき \\
    \hline
    第6問 & (1) & 軌跡は解答の図を参照 \\
    \cline{2-3}
    & (2) & 軌跡は解答の図を参照 \\
    \hline
  \end{tabular}
\end{table}

\paragraph{試験形態}
6問．試験時間は要確認．

\end{document}
